% YAAC Another Awesome CV LaTeX Template
%
% This template has been downloaded from:
% https://github.com/darwiin/yaac-another-awesome-cv
%
% Author:
% Christophe Roger
%
% Template license:
% CC BY-SA 4.0 (https://creativecommons.org/licenses/by-sa/4.0/)
%Section: Work Experience at the top
\definecolor{basecolor}{HTML}{251a85}       %BLUE

\sectionTitle{Experience}{}
%\renewcommand{\labelitemi}{$\bullet$}
\begin{experiences}

    \experience
    {Present}   {Institute of Foundational Models - MBZUAI [\href{https://ifm.ai/}{\textcolor{basecolor}{\faGlobe}}] | ML Research}{Rochester}{USA}
    {March 2026}
                    {
                    AI Research Intern | Advisor:\textit{\href{https://moonfolk.github.io/}{\textcolor{basecolor}{ Dr. Mikhail Yurochkin }}}
                      \begin{itemize}
                        \item Building novel reasoning datasets from web for the RL environments and SFT; synthetic data generation pipelines for improving the \textbf{K2 Think V2} capabilities on STEM reasoning tasks.
                      \end{itemize}
                    }{}
  \experience
    {February 2026}   {Rochester Institute of Technology [\href{https://www.rit.edu/}{\textcolor{basecolor}{\faGlobe}}] | ML Research}{Rochester}{USA}
    {September 2024}
                    {
                    Graduate Research Assistant | Advisor:\textit{\href{https://www.cs.cmu.edu/~akhudabu/}{\textcolor{basecolor}{Prof. Ashique KhudaBhuksh}}}
                      \begin{itemize}
                        \item Sole-authored publication at \textbf{NeurIPS 2025 (FM4LS)} on \textbf{Thin Bridges for Drug Text Alignment}, achieving \textbf{Recall@1 = 0.762} on \textbf{ChEMBL} and 3× scaffold-split gain over frozen baselines.
                        \item Co-authored publication in \textbf{INFORMS Journal on Computing}, developing a multimodal machine-comprehension framework evaluating 48{,}956 YouTube Kids videos; achieved \textbf{77.5\% accuracy} on ScienceQA and \textbf{F1 = 0.82} on the CVQA dataset.
                        \item Developing \textbf{diffusion-based brain decoding pipelines} in compute efficient way (\textit{fMRI → image reconstruction}) for multimodal perception tasks using \textbf{Stable Diffusion} and \textbf{CLIP} embeddings.
                      \end{itemize}
                    }{}

  \emptySeparator
    \experience
    {August 2024}   {Indian School of Business [\href{https://www.isb.edu/en.html}{\textcolor{basecolor}{\faGlobe}}] | Pre-doctorial Program}{Hyderabad}{India}
    {December 2020} {
                    Research Assistant | Advisor:\textit{\href{https://sumeetkumar.in/}{\textcolor{basecolor}{Prof. Sumeet Kumar}}
                    }
                      \begin{itemize}
                        \item Worked on \textbf{AI for Social Good} and \textbf{AI Safety} problems on social media platforms by processing large \textbf{multimodal datasets (YouTube, Amazon, Twitter)} at \textbf{TeraByte scale}, leading to publications at \textbf{AAAI, ASONAM, INFORMS, and ACL}.
                        \item Applied \textbf{PyTorch}, \textbf{OpenCV}, and \textbf{CUDA} to accelerate large-scale computer vision pipelines, optimizing training and inference speed for terabyte-scale datasets.
                      \end{itemize}
                    }{}
  \emptySeparator
    \experience
    {November 2020}   {Spacept [\href{https://spacept.com/}{\textcolor{basecolor}{\faGlobe}}] | Core ML Team}{Stockholm (Remote)}{Sweden}
    {August 2020} {
                    Research Intern | Advisor:\textit{\href{https://iliev.us/}{\textcolor{basecolor}{Sergiu Iliev ( Founder)}}
                    }
                      \begin{itemize}
                       \item Developed a \textbf{custom DL model} for \textbf{oil spill classification} using \textbf{Google Earth Engine} and Sentinel satellite data of the \textbf{Mauritius oil spill (July 2020)}, engineering scalable pipelines and augmentations that improved accuracy from \textbf{65\% to 78\%}, demonstrating transferability to real-world autonomous sensing.
                      \end{itemize}
                    }{}
  \emptySeparator
    \experience
    {March 2020}   {Indian Institute of Technology Madras [\href{https://www.iitm.ac.in/}{\textcolor{basecolor}{\faGlobe}}] | Neuromotive Team}{Chennai}{India}
    {December 2019} {
                    Research Intern | Advisor:\textit{\href{https://nvsteam.github.io/author/srinivasa-chakravarthy.-v/}{\textcolor{basecolor}{ Prof. Srinivasa Chakravarthy}}
                    }
                      \begin{itemize}
                        \item Implemented \textbf{Elman-Jordan-based CT reconstruction} halving projections (360 → 180) for \textbf{efficient imaging AI}, techniques extendable to \textbf{on-device simulation}.
                      \end{itemize}
                    }{}
  \emptySeparator
    \experience
    {July 2019}   {Indian Institute of Management Bangalore [\href{https://www.iimb.ac.in/home}{\textcolor{basecolor}{\faGlobe}}] | Research and Development Team}{Bangalore}{India}
    {May 2019} {
                    Research Intern | Advisor:\textit{\href{https://www.iimb.ac.in/user/138/trilochan-sastry}{\textcolor{basecolor}{ Prof.  Trilochan Sastry}}
                    }
                      \begin{itemize}
                        \item Analyzed the real-time data sales from \href{https://farmveda.in/?srsltid=AfmBOooAQs3NRcPk3ubLxGLKAerFbtsEoWFYjBDqH3DeC9mPPuLLVehq}{\textbf{Farmveda}} and predicted future sales, even further improving the sales by being actively involved in the Digital Marketing team.
                      \end{itemize}
                    }{}
    
  % \consultantexperience
  % {Janvier 2010}    {Ingénieur Consultant}{Altran Technologies}{France}
  % {Décembre 2007}   {IT Specialist}{IBM, Sensor Solutions Center of Excellence}
  %                   {
  %                     \begin{itemize}
  %                       \item Mise en place du suivi et du contrôle des commandes et approvisionnements à l'aide de la RFID
  %                       \item Projet de suivi et authentification de containers (\href{http://www.container-centralen.co.uk/rfid/history.aspx}{description}) : conception et développement
  %                       \item Amélioration d'une solution de contrôle des interventions dans un centre de données (RFID)
  %                       \item Solution de lutte contre la contrefaçon pour un producteur de spiritueux (RFID)           
  %                       \item Etude du protocole ONS : Analyse, \emph{POC}, documentation et présentation technique     
  %                       \item Maintenance corrective et évolutions d'une plateforme M2M (basée sur Websphere Portal)    
  %                     \end{itemize}
  %                   }
  %                   {
  %                     DB2, Eclipse, Infosphere Traceability Server, Lotus Expeditor, Eclipse, 
  %                     Rational Software Architect, IBM Premises Server, Maximo Asset Management for IT, RFIDIC (Infosphere Traceability Server - EPCIS)
  %                   }

\end{experiences}
